\subsection{Coefficient recovery across internal multiplicities}
\label{sec:coefficient-recovery}
The coefficient inverse, the extraction of scalar roots, and the
finite quotient of the target are separate operations. The first
continues analytically through an internal root multiplicity.

\begin{proposition}[Ambient coefficient recovery]\label{prop:coefficient-recovery}
Let $\theta_0\in\cK_d$ have strictly positive full-rank channels,
strict weights, $\tau(F(\theta_0))>0$, and pairwise distinct complete
monic component polynomials. Roots may repeat, both within and between
components. In monic-coefficient coordinates there is an ambient
real-analytic map $R$ near $P=F(\theta_0)$ with the following properties.
Every sufficiently nearby closed-model datum has precisely the
component permutations of $R(Y)$ as its coefficient representatives;
$R(P)=\theta_0$ for the selected labelling. Moreover, after removing
only the real-rootedness restriction in a neighbourhood of $\theta_0$,
$R\circ F$ is the identity on an open coefficient chart. In particular
$J_c=DF(\theta_0)$ is injective on that ambient chart. No manifold
assertion is made about its real-rooted subset at a multiple root.
\end{proposition}
\begin{proof}
The coefficient coordinates consist of the stochastic-column and
weight coordinates of Section~\ref{sec:local-condition} and the $kd$
nonleading polynomial coefficients. Positivity of the finitely many
clock values, strict stochastic inequalities, full channel ranks,
$\tau>0$, and distinct complete coefficient tuples persist on an open
set of these coordinates. This set does not require real-rootedness.

For an ambient observation $Y$, recover the monic normalizer by
$(\cN_Y^{\mathsf T}\cN_Y)^{-1}\cN_Y^{\mathsf T}b_Y$ and the numerator
by interpolation. Select the same independent leading-coefficient
rows and columns as at $P$, and form their analytic orthonormal frames
by positive definite Gram square roots. The reduced leading matrix is
invertible. A fixed real linear combination of its normalized
coefficient matrices has distinct eigenvalues at $P$, because the
complete polynomial tuples are distinct. Its disjoint Riesz contours
therefore give analytic rank-one projectors near $P$.
The traces against these projectors recover every component
coefficient. Formula~\eqref{eq:rank-one-recovery}, followed by total,
row, and column sums, recovers the weight and stochastic columns.
All these operations take place before scalar root extraction.

For data in the enlarged coefficient chart, the normalized
coefficient matrices still have their exact simultaneous diagonal
realization. The recovered projectors consequently give the original
coefficients and factors, so $R\circ F=\mathrm{id}$. Differentiating
proves the assertion about $J_c$. For a nearby closed-model datum,
uniqueness of the normalizer forces the same numerator; its rank-$k$
leading coefficient forces full-rank competing channels. The $k$
distinct rank-one joint projectors then force exactly one competing
component per tuple and the same row, column, and total sums. Hence
every competing coefficient representation is one of the finitely
many component permutations of $R(Y)$. Choose disjoint neighbourhoods
of these permutations and shrink the observation neighbourhood.
This proves the assertion for all representatives, not just for a
selected parametrization. Restriction of the ambient analytic map to
the real-rooted coefficient set completes the proof.
\end{proof}
